% ============================================================
% CONCLUSION
% ============================================================
\section{Conclusion}
\label{sec:conclusion}

This report presented a navigation system that takes a ground robot
through five random waypoints in a walled arena containing obstacles it
can only discover with its own sensors. The mission is planned on the
known walls: ant colony optimisation orders the waypoints by driving
distance, A* routes each leg on a map inflated for clearance, and
particle swarm optimisation smooths each leg towards the vehicle's
turning limit. During the mission, detected obstacles are remembered
and the remaining path is re-planned around them, with a reactive
potential field as a close-range backup. The Part~A estimator,
controller and waypoint optimiser were reused, and guidance was
extended to follow a path that can change during the mission.

Over $25$ random missions the vehicle captured every waypoint, passing
each within $0.23$~m, and, by the simulator's noise-free pose, never
came closer than $0.58$~m to an obstacle or $0.65$~m to a wall, a
margin of at least $0.18$~m over the required $0.40$~m. Re-planning is
what kept the obstacle clearance: with it switched off, six of eight
test missions breached it. Logging the noise-free pose also allowed the
sensor noise to be measured rather than assumed. It is zero-mean and
Gaussian with a standard deviation of $0.10$~m and $2^{\circ}$, and
recalibrating the Kalman filter to it made the filter consistent
without loss of accuracy. A mission took $123$~s on average, and of
the three speed limits tested, $1.0$~m/s was the fastest that kept the
clearance.

The main limitation is the trade between speed and margin. At
$1.5$~m/s the mission was $23$\% faster but one test mission breached
the clearance, because the vehicle cannot follow curves planned for a
$1.0$~m turn radius at that speed, and its own estimate can be over
$0.1$~m from the truth near an obstacle. Planning for a turn radius
that grows with the speed limit, and margins that grow with the
estimator's uncertainty, would allow a faster mission without losing
the margin. Two further limits apply outside simulation: obstacles are
assumed static, and each re-plan takes up to $0.65$~s that the
simulation does not charge for. An incremental planner such as D* Lite
\cite{koenig2002}, which repairs the previous search instead of
repeating it, would address the second.
